\subsection{实滤过}

定义 1. 设 G 是一个群。G 上的实滤过是指 G 的子群的一个族 \((\mathrm{G}_{\alpha})_{\alpha \in \mathbf{R}}\) ，使得

\[
\mathrm{G} _ {\alpha} = \bigcap_ {\beta <   \alpha} \mathrm{G} _ {\beta} \text {   for   all   } \alpha \in \mathbf {R}.\tag{1}
\]

公式 (1) 蕴涵 \(G_{\alpha} \subset G_{\beta}\) （对 \(\beta < \alpha\) ），因而族 \((\mathrm{G}_{\alpha})\) 是递减的。滤过 \((\mathrm{G}_{\alpha})\) 称为分离的，如果 \(\bigcap_{\alpha} G_{\alpha}\) 化为单位元素；称为穷竭的，如果 \(G = \bigcup_{\alpha} G_{\alpha}\) 。

注. 设 \((\mathrm{G}_n)_{n\in \mathbf{Z}}\) 是 G 的子群的一个递减序列。它是《交换代数》第 III 章 § 2 no. 1 定义 1 意义下的递减滤过。对每个整数 \(n\) 以及每个 \(\alpha\) （取自区间 \(]n - 1,n]\) ，即 \(\mathbf{R}\) 中的一个区间），我们记 \(\mathrm{H}_{\alpha} = \mathrm{G}_{n}\) ，特别地 \(\mathrm{H}_n = \mathrm{G}_n\) 。立即可见我们这样就得到 G 上的一个实滤过 \((\mathrm{H}_{\alpha})_{\alpha \in \mathbf{R}}\) ；这样的滤过将称为整滤过。因此《交换代数》第 III 章 § 2 意义下的递减滤过可以与整滤过等同。

设 A 是一个代数；A 的加法群上的实滤过 \((\mathrm{A}_{\alpha})\) 称为与该代数结构相容，如果 \(\mathrm{A}_{\alpha}.\mathrm{A}_{\beta}\subset \mathrm{A}_{\alpha +\beta}\) （对 R 中的 \(\alpha ,\beta\) ），且 \(\mathrm{K.A}_{\alpha}\subset \mathrm{A}_{\alpha}\) （对 \(\alpha \in \mathbf{R}\) ）。若滤过是穷竭的，则 \((\mathrm{A}_{\alpha})\) 是 A 上与该代数结构相容的拓扑之下 0 的一个邻域基本系。设 B 是一个幺代数；B 的加法群上的实滤过 \((\mathrm{B}_{\alpha})\) 称为与该幺代数结构相容，如果它与该代数结构相容并且 \(1\in \mathrm{B}_0\) 。

